% AUTORIDAD DE FUENTE:
% expose_VIII_body.tex, líneas 2243--2261; índices de escaneo 313--314;
% folios de fuente 302--303. No hay diferencias materiales entre capas.

\medskip
\noindent\underline{Observación 7.2.} Sea $P$ un esquema en grupos liso
sobre un esquema $S$ conexo, regular y localmente noetheriano, con
punto genérico $\eta$. Cabe estudiar el funtor $E\mapsto E_\eta$ de la
categoría de extensiones de $P$ por $\underline G_{mS}$ a la categoría
de extensiones de $P_\eta$ por $\underline G_{m\eta}$. El problema se
reduce fácilmente al tratado en 7.1(a). Para simplificar, supongamos que
$P$ es de tipo finito sobre $S$. Se deduce fácilmente de 7.1 que este
funtor es plenamente fiel.\footnote{Para esto basta suponer que $S$ es
normal, en vez de regular.} De aquí se deduce que, si $E_\eta$ es una
extensión de $P_\eta$ por $\underline G_{m\eta}$, entonces es isomorfa a
la fibra genérica de una extensión $E$ de $P$ por $\underline G_{mS}$
si y solo si, para todo $s\in S$ tal que
\[
 \dim\operatorname{Spec}(\underline O_{S,s})=1,
\]
el problema análogo de extensión obtenido por el cambio de base
$\operatorname{Spec}(\underline O_{S,s})\to S$ tiene solución. Por
7.1(a), esto se expresa mediante la anulación de cierto homomorfismo
\begin{equation}\tag{7.2.1}
 d(E_\eta,s):P_s/P_s^{\circ}
 \longrightarrow(\mathbb Q/\mathbb Z)_{k(s)}.
\end{equation}

%%%%%%%%%%  scan idx 314  |  folio 303  %%%%%%%%%%

La afirmación precedente se reduce manifiestamente a la siguiente.
Sea $U$ un abierto de $S$ cuyo complemento $T$ satisface
\[
 \operatorname{codim}(T,S)\geq2.
\]
Entonces el funtor $E\mapsto E|U$ de la categoría de extensiones de
$P$ por $\underline G_{mS}$ a la categoría de extensiones de $P_U$ por
$\underline G_{mU}$ es una equivalencia de categorías. En términos de
la descripción VII 1.1.6, esta afirmación es esencialmente trivial:
para todo esquema liso $X$ sobre $S$ (como $P$, $P\times P$,
$P\times P\times P$, etc.), $X$ es localmente factorial y
$X_T=X\setminus X_U$ tiene codimensión al menos $2$ en $X$. Por tanto,
el funtor $F\mapsto F|X_U$ de la categoría de torsors bajo
$\underline G_{mX}$ a la categoría de torsors bajo
$\underline G_{mX_U}$ es una equivalencia de categorías.

Hay resultados completamente análogos para el problema de extender
biextensiones; dejamos su formulación al lector.
